\documentclass[10pt,a4paper]{article}
\usepackage[T1]{fontenc}
\usepackage[utf8]{inputenc}
\usepackage{lmodern,microtype,amsmath,amssymb,booktabs}
\usepackage[margin=19mm,top=17mm,bottom=18mm]{geometry}
\usepackage{xcolor,hyperref}
\definecolor{link}{HTML}{174E70}
\hypersetup{colorlinks=true,linkcolor=link,urlcolor=link,
  pdftitle={Hecke calculations for the conditional vanishing of K4k(Z) through degree 104},
  pdfauthor={Leo Pintar},pdfdisplaydoctitle=true}
\urlstyle{same}
\setlength{\parindent}{0pt}
\setlength{\parskip}{3.5pt}
\setlength{\abovedisplayskip}{5pt}
\setlength{\belowdisplayskip}{5pt}
\newcommand{\Z}{\mathbb Z}
\newcommand{\Q}{\mathbb Q}
\newcommand{\E}[1]{\Q(\sqrt{-#1})}
\newcommand{\tr}{\operatorname{tr}}
\newcommand{\num}{\operatorname{num}}
\newcommand{\subhead}[1]{\par\smallskip{\large\bfseries #1}\par}
\newcommand{\repourl}{https://github.com/leopintar05/hecke-k-theory-104/releases/tag/v1.1}
\title{Hecke calculations for the conditional vanishing\\[2pt]of \(K_{4k}(\Z)\) through degree 104}
\author{Leo Pintar}
\date{Revised 21 November 2026}
\makeatletter
\renewcommand{\maketitle}{%
\begin{center}
{\LARGE\bfseries\@title\par}
\vspace{5pt}
{\large\@author}\quad École Polytechnique\\[2pt]
\@date
\end{center}}
\makeatother
\renewenvironment{abstract}{\par\noindent\textbf{Abstract.}\enspace\ignorespaces}{\par}
\begin{document}
\maketitle
\begin{abstract}
In degree 104, two rank-324 Hecke operators over \(\E{3}\), together with classical tests, leave four candidate primes. An integer evaluation matrix of rank 2,925 constructed using \(\E{7}\) has determinant coprime to their product. Together with the certificates for earlier degrees, these computations imply \(K_{4k}(\Z)=0\) for \(1\le k\le26\), assuming the strict Selmer comparison, coverage of the admissible representations, and a simultaneous residual trace condition. The supplement contains the matrices or their generators, exact residues, and replay code.
\end{abstract}

\subhead{1. The degree-104 certificate}
Put \(e=52\), \(E_d=\E{d}\) for \(d=3,7\), and
\(b_j=|\num(B_{j,\eta}/j)|\), where \(\eta\) is the quadratic character of conductor 3 and the numerator is taken in lowest terms. Define \(p_1,p_2,p_3,p_4\) by
\[
\begin{gathered}
 p_1=43\,279\,339\,111,\qquad p_2=490\,440\,141\,803,\\
 p_3=63\,421\,406\,191\,963\,401\,438\,831\,238\,997,\qquad p_4=b_{105}/2162.
\end{gathered}
\]
Primality tests and exact arithmetic give \(b_{53}=2p_1p_2p_3\) and \(b_{105}=2\cdot23\cdot47\,p_4\). Write \(M=p_1p_2p_3p_4\).

\textit{First field.} On the rank-324 negative subspace for \(E_3\), shift the inert-2 operator by \(3\cdot4^{51}-(4^{104}+4^{52}+1)\) to obtain \(X\). Shift each split-7 operator by \(-(7^{104}+7^{52}+1)\) and take their sum \(Z\). Reconstruction from the neighbor sums gives
\[
 \gcd(\det X,\det Z)=3^{1578}5^{33}11^{31}19^{18}23^6\cdot29\cdot79\,p_1p_2p_3 .
\]
For a non-self-dual witness the difference of conjugate eigenvectors forces \(\ell\mid\det X\) and \(\ell\mid\det Z\). The matrices have a common kernel line modulo each of \(p_1,p_2,p_3\). The classical branch has weight-105 trace determinants \(C_q\) with \(4\gcd(C_2,C_7)=3^{166}5^{12}11^4 19^6 b_{105}^2\), leaving only \(p_4\) above the cutoff.

\textit{Self-dual branch over \(E_3\).} By coverage the witness is the adjoint of \(r\) attached to an eigenform in \(S_{53}(\Gamma_0(3),\eta)\). For rational Frobenius \(F_q\), write \(a_q=\tr r(F_q)\) and \(\det r(F_q)=\eta(q)q^{52}\). The places over 2 and 7 have norms \(Q=4,7\) and Frobenius \(g_2=F_2^2,\ g_7=F_7\), respectively. Set \(\Delta_q=\tr\operatorname{Ad}^0(r(g_q))-(1+Q^{52}+Q^{-52})\). Since \(\tr\operatorname{Ad}^0(A)=(\tr A)^2/\det A-1\),
\[
  2^{104}\Delta_2=(a_2^2+2^{53})^2-(1+2^{104})^2,\qquad
  7^{52}\Delta_7=a_7^2-(1+7^{52})^2 .
\]
Let \(H_q\) be the right-hand polynomial evaluated at \(T_q\). A witness at both places forces \(\ell\) to divide both determinants. Fourier expansions through order 112 give two \(16\times16\) matrices with
\[
 \gcd(\det H_2,\det H_7)=3^{76}5^{12}11^2 19^2(p_1p_2p_3)^2.
\]
Thus every possible prime \(\ell\ge2^{31}\) in the three branches divides \(M\).

\textit{Second field.} At both places above 2 in \(E_7\), the scaled target is \(u=2^{104}+2^{52}+1\). The level-seven tests in weights 53 and 105 exclude the self-dual and classical branches at all four primes. In the remaining case conjugation \(J\) exchanges distinct eigensystems. For an eigenvector \(v\), the difference \(v-Jv\ne0\) is a negative eigenvector for the sum of the two operators. Its eigenvalue reduces to \(t=2u\) under the simultaneous trace hypothesis.

The two lattice classes give directed operators with blocks \(\left(\begin{smallmatrix}A&B\\D&0\end{smallmatrix}\right)\) and \(\left(\begin{smallmatrix}A'&B'\\D'&0\end{smallmatrix}\right)\). As \(\gcd(t,M)=1\), elimination of the second class gives the operator on the standard negative subspace
\[
 L=t^2I-t(A+A')-(B+B')(D+D').
\]
Its dimension is \(\binom{27}{3}=2925\). Let \(F_j=P_j-JP_j\) be the signed orbit basis of bidegree \((51,51)\) modulo \(x\cdot y\). At the integral incidence pairs in the supplement (seed 128711) evaluate \(V_{ij}=F_j(x_i,y_i)\) and \(W_{ij}=(LF_j)(x_i,y_i)\). Then \(W=V[L]\) over \(\Q\). The scaled correspondences preserve a lattice at each \(p_i\), so the proposed eigenvalue forces positive \(p_i\)-adic valuation of \(\det L\) and hence \(p_i\mid\det W\).

The residues \(r_i=\det W\bmod p_i\) are all nonzero, so \(\gcd(\det W,M)=1\). The prime bit lengths are \(36,39,96,432\).
\begin{center}
\begin{tabular}{@{}cl@{\qquad}l@{}}\toprule
\(i\)& \(r_i\) in decimal & \(r_4\) in decimal (concatenate the three lines)\\\midrule
1 & \texttt{32189186543} & \texttt{385227761542015351632264879402615235505177920}\\
2 & \texttt{445745447561} & \texttt{099210521658726589115415728765387796061514558}\\
3 & \texttt{22450212846525469764552672927} & \texttt{7434547466097567956069671093168965749031}\\\bottomrule
\end{tabular}
\end{center}
\newpage
\subhead{2. Arithmetic hypotheses and the range through degree 104}
Fix \(e=2k\), \(1\le k\le26\), and a prime \(\ell\ge2^{31}\). Write \(\omega\) for the residual cyclotomic character and \(S_e\) for the subgroup of \(H^1(\Q,(\Q_\ell/\Z_\ell)(-e))\) consisting of classes unramified everywhere and trivial at 2 and \(\ell\). The deduction uses the following three inputs in the fields and at the primes tested.

\textit{Strict Selmer comparison.} A nonzero \(K_{2e}(\Z)/\ell\) gives a nonzero class in \(S_e\). Assume that the length in Alpöge's \texttt{OneCharU\_s}, for \(N=2\), trivial character and coefficients \(\Z_\ell\), measures this strict group. The required implication, at a fixed coefficient embedding, is
\[
 S_e\ne0\ \Longrightarrow\
 \forall \sigma\in G_{E_d}\ \exists\,\rho_\sigma\ \text{admissible}:\quad
 \tr\bar\rho_\sigma(\sigma)=(1+\omega^e+\omega^{-e})(\sigma).
\]
If all differences at one \(\sigma\) were units, the valuation sum in the bound would be zero. The source defines the bound in \href{https://github.com/CocoSgt/riemann-hypothesis-lean/blob/7cbf7a0c32fca1badf05931879a37c3f8a2b03b4/src/Development/AutomorphicArithmetic/AutomorphicVarieties/ModularityArgument/ThresholdArguments/DenseThresholdCases/NumberU3US.lean#L46-L59}{\texttt{NumberU3US.lean}} and constructs it inside \href{https://github.com/CocoSgt/riemann-hypothesis-lean/blob/7cbf7a0c32fca1badf05931879a37c3f8a2b03b4/src/Development/AutomorphicArithmetic/AutomorphicVarieties/SmallClassNumbers/CorGKOf.lean#L10-L25}{\texttt{corGK\_of}}.

\textit{Coverage.} Assume the systems admitted by \texttt{Adm3U\_s} and \texttt{Adm2A} occur with the tested weights, levels, twists and Frobenius normalizations. This includes the transfer to the compact unitary space of weight \((e-1,0,1-e)\), the self-dual reduction and the classical spaces. For the two-class model see \href{https://arxiv.org/abs/0801.3176}{Loeffler, \emph{Explicit calculations of automorphic forms for definite unitary groups}}.

\textit{Simultaneous residual traces.} For each field separately, assume that \(S_e\ne0\) supplies one admitted system at one coefficient prime realizing the target traces at all places compared together. The pointwise implication permits \(\rho_\sigma\) to vary. A proposed interpolation on ordinary inertia uses a bound of \(s=2e^3+6(4e+3)\) candidate systems, giving \(2es=29\,378\,128<\ell-1\) at \(e=52\).

\textbf{Conditional conclusion.}
Under these inputs \(K_{4k}(\Z)=0\) for \(1\le k\le26\).
For odd \(\ell\), the cyclotomic comparison detects the quotient by \(\ell\) in the even \(\omega^{-e}\)-component of the cyclotomic class group. Vandiver's conjecture has been checked throughout the required range \(\ell<2^{31}\). The known 2-primary calculation treats \(\ell=2\). See \href{https://arxiv.org/html/1703.00108v3#S5}{Jordan et al., Theorem 5.1}, \href{https://arxiv.org/abs/1605.02398}{Hart--Harvey--Ong, \emph{Irregular primes to two billion}}, and \href{https://sites.math.rutgers.edu/~weibel/Kbook/Kbook.VI.pdf}{Weibel, \emph{The \(K\)-book}, VI.10}.

For \(e=2,4,\ldots,50\), let \(C_q,U_q\) be the classical and compact comparison determinants over \(E_3\) at \(q=2,7\). Their archived integer values satisfy, for every prime \(\ell>337\),
\[
 v_\ell\!\left(\gcd(C_2U_2,C_7U_7)\right)
 =v_\ell(b_{e+1})+3v_\ell(b_{2e+1}).
\]
Thus all large primes requiring a second-field test divide \(b_{e+1}b_{2e+1}\). In degrees \(52,56,60,64,68\) the first field leaves respectively \(4,2,4,4,3\) primes. The second-field determinants exclude them. In degree 72 the test works modulo a product containing all survivors. The supplement includes these cases and the earlier degrees.

For degrees \(76,80,\ldots,100\), removing certified small factors from \(b_{e+1}b_{2e+1}\) gives a test modulus \(M_e\). The block elimination of Section 1 gives a signed evaluation matrix \(W_e\). The recorded results are
\begin{center}
\begin{tabular}{@{}lrrrrrrr@{}}\toprule
Degree & 76 & 80 & 84 & 88 & 92 & 96 & 100\\\midrule
Matrix dimension & 1140 & 1330 & 1540 & 1771 & 2024 & 2300 & 2600\\
Bit length of \(M_e\) & 368 & 419 & 426 & 475 & 489 & 515 & 527\\
\(\gcd(\det W_e,M_e)\) & 1 & 1 & 1 & 1 & 1 & 1 & 1\\\bottomrule
\end{tabular}
\end{center}
Section 1 supplies the four-prime exclusion for degree 104. These groups are finite, so exclusion at every prime gives the conditional conclusion.

\textbf{Prior work and reproducibility.}
Alpöge suggested this application in \href{https://alpo.ge/remark.pdf#page=4}{Step 17, footnote 6}. A \href{https://github.com/KokunoYumeto/grothendieck-montpellier-archive/blob/daddc785c90b0447f448ab372b1014efa2afb208/research/grh_20261120/intakes/K17/INTAKE_K17_20261121_part4.md#L69-L101}{21 November report in the Grothendieck--Montpellier archive} gives a conditional calculation through degree 48. The \href{\repourl}{v1.1 supplement} covers the fourteen additional degrees \(52,56,\ldots,104\). Run \texttt{python check\_adjoint104.py} for the two weight-53 traces. With the listed dependencies, \texttt{python reproduce.py} also rebuilds the two rank-324 matrices and the rank-2,925 matrix. It reduces the latter from the 602-bit modulus \(M\) to each prime field and independently checks three complete rows by evaluating all 40 oriented contributions.

{\small\textit{Contributions.} Leo Pintar initiated and led this investigation into the \(K\)-theoretic consequences of Alpöge’s Step 17. OpenAI Codex contributed to the mathematical arguments, wrote and ran the verification code, and drafted the manuscript.}
\end{document}
